%% content/shared.tex -- text shared by the slide deck (main.tex) and the
%% one-page handout (handout.tex), so the two can never disagree.
%%
%% EDIT BEFORE PRESENTING: the three placeholders below.
\newcommand{\Presenter}{[Presenter name]}
\newcommand{\ContactLine}{[email or phone]}
\newcommand{\StartDate}{January 2027}  % sample calendar in each Teacher Guide

%% ---------------------------------------------------------------- course facts
%% Source of truth: grade-4-6/SYLLABUS.md, docs/DESIGN.md, docs/ASSUMPTIONS.md.
\newcommand{\CourseName}{Kids Coding}
\newcommand{\LevelWeeks}{10--12 weeks}   % per level; holiday weeks without a session do not count
\newcommand{\TotalWeeks}{30--36 weeks}   % all three levels
\newcommand{\CourseShape}{3 levels \textperiodcentered{} \LevelWeeks{} each}
\newcommand{\SessionLength}{one 60-minute session a week}
\newcommand{\GroupSize}{up to 4 children, working in 2 pairs}
\newcommand{\Homework}{15--20 minutes a week, optional but encouraged}
\newcommand{\Equipment}{a Mac, the free Scratch app and Thonny; no accounts, no internet during lessons}

%% Level summaries: \LevelXName, \LevelXTools, \LevelXLearn, \LevelXProject, \LevelXProjectLine
\newcommand{\LevelOneName}{Introductory}
\newcommand{\LevelOneTools}{Scratch blocks, then Python turtle}
\newcommand{\LevelOneLearn}{steps in order, loops, events and choices, variables, first Python}
\newcommand{\LevelOneProject}{My Interactive Art Show}
\newcommand{\LevelOneProjectLine}{a program that asks the viewer questions and draws a picture that changes with their answers}

\newcommand{\LevelTwoName}{Intermediate}
\newcommand{\LevelTwoTools}{Python}
\newcommand{\LevelTwoLearn}{functions, lists and dictionaries, files, testing, keyboard-driven games}
\newcommand{\LevelTwoProject}{Arcade Game or Art Studio}
\newcommand{\LevelTwoProjectLine}{an arcade game with score and high scores, or an art studio that saves designs, with the child's own tests}

\newcommand{\LevelThreeName}{Advanced}
\newcommand{\LevelThreeTools}{Python projects + git}
\newcommand{\LevelThreeLearn}{objects and classes, search and sort, recursion, data, how simple AI learns}
\newcommand{\LevelThreeProject}{Build Something Big}
\newcommand{\LevelThreeProjectLine}{a grid game, art app or data/AI app with classes, an algorithm, tests and a git history}

%% ---------------------------------------------------------------- how parents help
\newcommand{\HelpItems}{%
  \item \textbf{Ask for a demo:} \enquote{Show me what you made.}
  \item \textbf{Ask about the bug,} not the score: \enquote{What went wrong? How did you find it?}
  \item \textbf{Praise strategies and effort,} not \enquote{being a natural.}
  \item \textbf{Protect 15 minutes} a week for homework, at a regular time.
  \item \textbf{Encourage everyone equally:} girls hear less encouragement than boys.
  \item \textbf{Come to the showcase} at the end of each level.}

%% ---------------------------------------------------------------- safety promises
\newcommand{\SafetyItems}{%
  \item \textbf{No accounts, no logins, no email addresses.}
  \item \textbf{Nicknames only;} programs never ask for personal information.
  \item \textbf{Nothing is uploaded;} projects stay on the computers in the room.
  \item \textbf{Screen breaks every session,} plus unplugged activities.
  \item \textbf{Color-blind-safe materials} and a dyslexia-friendly edition.}
